\section{The differential equation and its local solutions}\label{sec:ode}

\subsection{Coefficient proof of the equation}

Let
\[
 \mathcal L_{a,b,c}y
 =z(1-z)y''+[c-(a+b+1)z]y'-aby.
\]

\begin{theorem}[Hypergeometric differential equation]\label{thm:ode}
The formal Gauss series satisfies
\begin{equation}\label{eq:hg-ode}
 \mathcal L_{a,b,c}\F a b c z=0.
\end{equation}
On every strict subdisk where the series is effectively convergent, the same
identity holds between named analytic functions.
\end{theorem}

\begin{proof}
For \(y=\sum_{n\ge0}A_nz^n\), the coefficient of \(z^n\) in
\(\mathcal L_{a,b,c}y\) is
\[
 (n+1)(n+c)A_{n+1}
 -[n^2+(a+b)n+ab]A_n
 =(n+1)(n+c)A_{n+1}-(n+a)(n+b)A_n.
\]
It vanishes by \cref{lem:ratio}.  This is a single rational identity for
each \(n\).  Analytic equality follows from
\cref{thm:analytic-derivative}: all three named series have explicit common
tail moduli, so coefficientwise cancellation passes to their sums.
\end{proof}

\begin{corollary}[Normalized regular solution is unique]\label{cor:ode-unique}
Suppose \(y=\sum B_nz^n\), \(B_0=1\), and
\(\mathcal L_{a,b,c}y=0\).  If \(n+c\) is apart from zero for every observed
\(n\), then \(y=\F a b c z\) formally.
\end{corollary}

\begin{proof}
The ODE coefficient gives exactly the recurrence in
\cref{thm:formal-uniqueness}.
\end{proof}

\subsection{Frobenius calculation at the three singular points}

The following calculations are algebra in a ring with a formal monomial
\(z^s\) and the derivative rule
\(D(z^s)=sz^{s-1}\).  Analytically, the monomial is interpreted using a
certified logarithm branch.

\begin{proposition}[Solutions at zero]\label{prop:zero-solutions}
The indicial polynomial at \(z=0\) is
\[
 I_0(s)=s(s+c-1).
\]
The two formal Frobenius expressions are
\begin{align}
 u_{0,1}(z)&=\F a b c z,\label{eq:u01}\\
 u_{0,2}(z)&=z^{1-c}
   \F{a+1-c}{b+1-c}{2-c}{z}.
   \label{eq:u02}
\end{align}
Under nonresonance and branch certificates they are named local analytic
solutions.
\end{proposition}

\begin{proof}
Substitute \(y=z^sv\).  The coefficient of the lowest power is
\(s(s+c-1)v(0)\).  For \(s=0\), the equation for \(v\) is the original one.
For \(s=1-c\), direct collection gives
\[
 \mathcal L_{a,b,c}(z^{1-c}v)
 =z^{1-c}\mathcal L_{a+1-c,b+1-c,2-c}v.
\]
The claimed solutions now follow from \cref{thm:ode}.  The displayed gauge
identity is checked after clearing the monomial by a finite polynomial
calculation in \(s\); substituting \(s=1-c\) finishes it.
\end{proof}

\begin{proposition}[Solutions at one]\label{prop:one-solutions}
With \(t=1-z\), the exponents at \(z=1\) are \(0\) and \(c-a-b\), and two
standard expressions are
\begin{align}
 u_{1,1}(z)&=\F a b{a+b+1-c}{1-z},\label{eq:u11}\\
 u_{1,2}(z)&=(1-z)^{c-a-b}
  \F{c-a}{c-b}{c+1-a-b}{1-z}.
  \label{eq:u12}
\end{align}
\end{proposition}

\begin{proof}
Replace \(z\) by \(1-t\) and use \(d/dz=-d/dt\).  After collecting terms,
the regular equation in \(t\) has parameters
\((a,b,a+b+1-c)\).  The second expression is obtained from the second
exponent exactly as in \cref{prop:zero-solutions}.  Every step is an identity
of rational differential operators.
\end{proof}

\begin{proposition}[Solutions at infinity]\label{prop:infty-solutions}
With \(w=z^{-1}\), the exponents at infinity are \(a,b\) in the convention
\(y\sim z^{-s}\).  Two expressions are
\begin{align}
 u_{\infty,a}(z)&=(-z)^{-a}
   \F a{a+1-c}{a+1-b}{z^{-1}},\label{eq:uinf-a}\\
 u_{\infty,b}(z)&=(-z)^{-b}
   \F{b+1-c}{b}{b+1-a}{z^{-1}}.
   \label{eq:uinf-b}
\end{align}
The choice of \((-z)^{-a}\) rather than \(z^{-a}\) is branch bookkeeping;
formally either monomial is a unit multiple of the other.
\end{proposition}

\begin{proof}
Use
\(D_z=-w^2D_w\) and
\(D_z^2=w^4D_w^2+2w^3D_w\), then set
\(y=z^{-s}v(w)\).  The lowest-order equation is
\((s-a)(s-b)=0\).  At \(s=a\) and \(s=b\), collection gives the two displayed
parameter triples.  A finite symbolic checker can verify the operator
identity after multiplying by a sufficient power of \(z\).
\end{proof}

\begin{remark}[Resonance]
If an exponent difference is an integer, the second local solution may
contain a logarithm.  The nonlogarithmic formulas can coincide or become
undefined because a denominator parameter is a nonpositive integer.  The
constructive response is not to pretend that the generic basis persists.
For every fixed truncation order, substitute
\[
 y=y_1\log z+z^s\sum_{n<N}d_nz^n+O(z^{s+N})
\]
into the differential equation.  The coefficients \(d_n\) satisfy a finite
inhomogeneous recurrence involving finite harmonic sums (equivalently,
digamma differences).  This recurrence is primitive recursive.  A full
uniform logarithmic solution is a \StatusB{} theorem once a joint coefficient
bound and convergence modulus are supplied.  The present package gives the
recurrence interface but does not claim a completed all-resonances analytic
library.
\end{remark}

\subsection{Wronskians and local independence}

\begin{lemma}[Wronskian equation]\label{lem:wronskian}
For two solutions of \eqref{eq:hg-ode}, their Wronskian satisfies
\[
 W'(z)=-\frac{c-(a+b+1)z}{z(1-z)}W(z).
\]
Equivalently
\[
 W(z)=Cz^{-c}(1-z)^{c-a-b-1}
\]
on a certified branch domain.
\end{lemma}

\begin{proof}
For an equation \(y''+P y'+Qy=0\), direct expansion gives
\(W'=-PW\).  Here
\[
 P(z)=\frac{c-(a+b+1)z}{z(1-z)}
      =\frac c z+\frac{c-a-b-1}{1-z}.
\]
The displayed solution follows by differentiating the right side.  Thus the
claim can be verified without invoking indefinite integration.  A nonzero
constant \(C\) is witnessed by the leading terms of a nonresonant Frobenius
pair.
\end{proof}

\subsection{The normal or Q-form}

\begin{proposition}[Elimination of the first derivative]\label{prop:q-form}
Set
\[
 y(z)=v(z)u(z),
 \qquad
 v(z)=z^{-c/2}(1-z)^{(c-a-b-1)/2}.
\]
Then \(u\) satisfies
\[
 u''+Q(z)u=0.
\]
After choosing the convention consistent with \(y=vu\), where
\[
 Q(z)=
 \frac{z^2[1-(a-b)^2]
       +z[2c(a+b-1)-4ab]+c(2-c)}
      {4z^2(1-z)^2}.
\]
\end{proposition}

\begin{proof}
Divide the original equation by \(z(1-z)\), obtaining
\(y''+Py'+Ry=0\).  For \(y=vu\), the coefficient of \(u'\) is
\(2v'/v+P\), which vanishes because
\[
 \frac{v'}v=-\frac{c}{2z}-\frac{1+a+b-c}{2(z-1)}=-\frac P2.
\]
The remaining coefficient is \(R-P'/2-P^2/4\).  Clearing the denominator
\(4z^2(1-z)^2\) and expanding yields the displayed numerator.  This is a
polynomial identity suitable for the repository's ring tactic.
\end{proof}

\subsection{Reduction of a three-singularity equation}

The classical slogan that every second-order linear equation with three
regular singular points is hypergeometric must be made data-sensitive.

\begin{theorem}[Certified Riemann-to-Gauss reduction]\label{thm:riemann-reduction}
Let
\[
 y''+P(x)y'+Q(x)y=0
\]
be a rational differential equation.  Supply:
\begin{enumerate}
\item three pairwise apart algebraic points \(x_0,x_1,x_\infty\), with a
certificate that the poles of \(P,Q\) occur only there;
\item regular-singularity bounds
\((x-x_i)P(x)\) and \((x-x_i)^2Q(x)\) at each finite point, and the analogous
bounds at infinity;
\item isolated roots of the three indicial polynomials;
\item the Fuchs relation between the six exponents.
\end{enumerate}
Then a primitive recursive symbolic procedure constructs a fractional-linear
coordinate \(z\) sending the three points to \(0,1,\infty\), a gauge
\(y=z^\alpha(1-z)^\beta u\), and parameters \(a,b,c\) for which \(u\)
satisfies the hypergeometric equation.  The correctness proof is a finite
rational-function identity.
\end{theorem}

\begin{proof}
Take
\[
 z=\frac{(x-x_0)(x_1-x_\infty)}
          {(x-x_\infty)(x_1-x_0)}.
\]
Apartness makes every denominator invertible.  The chain rule transforms the
equation to one with singularities at \(0,1,\infty\).  Choose the gauge
exponents to make one exponent at \(0\) and one at \(1\) equal to zero.  If
the remaining exponent differences are \(1-c\), \(c-a-b\), and \(a-b\),
solve the resulting linear equations for \(a,b,c\).  The Fuchs relation
ensures consistency.  After multiplication by a common polynomial
denominator, equality with \(\mathcal L_{a,b,c}\) is an exact coefficient
identity.  All root choices were supplied as isolated algebraic names, so no
unbounded selection is hidden in the construction.
\end{proof}
